
\documentclass{article}
\usepackage{colortbl}
\usepackage{makecell}
\usepackage{multirow}
\usepackage{supertabular}

\begin{document}

\newcounter{utterance}

\centering \large Interaction Transcript for game `function\_detective', experiment `two\_numbers\_passive\_labeled\_pairs', episode 297 with mistral{-}large{-}3{-}azure.
\vspace{24pt}

{ \footnotesize  \setcounter{utterance}{1}
\setlength{\tabcolsep}{0pt}
\begin{supertabular}{c@{$\;$}|p{.15\linewidth}@{}p{.15\linewidth}p{.15\linewidth}p{.15\linewidth}p{.15\linewidth}p{.15\linewidth}}
    \# & \multicolumn{2}{c}{Player} && \multicolumn{2}{c}{Game Master} \\
    \hline

    \theutterance \stepcounter{utterance}  
    & & & \multicolumn{4}{p{0.6\linewidth}}{
        \cellcolor[rgb]{0.9,0.9,0.9}{
            \makecell[{{p{\linewidth}}}]{
                \texttt{\tiny{[P1$\langle$GM]}}
                \texttt{I have a hidden Python function for you to figure out.} \\
\\ 
\texttt{Your job is to infer what it does and then submit a correct Python implementation.} \\
\\ 
\texttt{Function signature:} \\
\texttt{(a: int, b: int) {-}> int} \\
\\ 
\texttt{You have 15 turns total.} \\
\\ 
\texttt{I reveal candidate input{-}output pairs one at a time and tell you whether each one is correct.} \\
\\ 
\texttt{You must make your best SOLVE attempt no later than the second{-}to{-}last turn.} \\
\texttt{Do not wait until the very end.} \\
\\ 
\texttt{You can use these actions:} \\
\texttt{{-} JSON format described at the end of the prompt} \\
\texttt{{-} SOLVE: ```python ... ```} \\
\\ 
\texttt{Do not add import statements in your guess.} \\
\texttt{Anything you might need is already available, so just write the function body in hypothesis and return the answer.} \\
\\ 
\texttt{Preloaded examples:} \\
\texttt{1. (((1, 10), 3), True)} \\
\\ 
\texttt{Scoring:} \\
\texttt{{-} correctness is required for any non{-}zero score} \\
\texttt{{-} in interactive modes, solving with fewer information requests is better} \\
\\ 
\texttt{Input formatting reminders:} \\
\texttt{{-} strings must be quoted} \\
\texttt{{-} booleans must be True/False} \\
\texttt{{-} use commas between multiple values} \\
\texttt{{-} do not write any extra action tags in the same message} \\
\\ 
\texttt{Format example:} \\
\texttt{You write:} \\
\texttt{\{'current\_mode': <probing/confirming/don't know; choose whichever best describes your current state in the task>, 'input\_rationale': <brief description of why you chose this hypothesis>, 'hypothesis': <either the exact string "unknown" if you have no guess yet OR your current best guess as runnable Python matching the given signature, with a one{-}line docstring stating what it does. Never write placeholder code, if you cannot commit to a guess, use "unknown". Match the given signature.} \\
\texttt{e.g.} \\
\texttt{def f(x):} \\
\texttt{	"""Returns the absolute value of x."""} \\
\texttt{	return x if x >= 0 else {-}x} \\
\\ 
\texttt{Use "unknown" if you have no guess.>, 'confidence\_score': <score between 1{-}100. 100 means highest confidence in your hypothesis>\}} \\
\\ 
\texttt{Output only one valid json per turn.} \\
            }
        }
    }
     \\ \\

    \theutterance \stepcounter{utterance}  
    & \multicolumn{4}{p{0.6\linewidth}}{
        \cellcolor[rgb]{0.561,0.737,0.561}{
            \makecell[{{p{\linewidth}}}]{
                \texttt{\tiny{[P1$\rangle$GM]}}
                \texttt{```json} \\
\texttt{\{} \\
\texttt{  "current\_mode": "probing",} \\
\texttt{  "input\_rationale": "The preloaded example shows that f(1, 10) returns 3. This could be related to the number of integers between 1 and 10 that satisfy a certain property (e.g., prime numbers, odd numbers, etc.). The number 3 could correspond to the count of primes between 1 and 10 (2, 3, 5, 7), but 1 is not a prime, so it would be 4 primes. Alternatively, it could be the count of odd numbers (1, 3, 5, 7, 9), which is 5. Another possibility is the count of numbers with exactly one '1' in their binary representation (1, 2, 4, 8), which is 4. The number 3 doesn't match these, so it might be something else, like the count of numbers where the sum of digits is odd (1:1, 2:2, 3:3, 4:4, 5:5, 6:6, 7:7, 8:8, 9:9, 10:1) → 1,3,5,7,9,10 → 6 numbers. Not matching. Alternatively, it could be the count of numbers where the number of set bits is odd} \\
            }
        }
    }
    & & \\ \\

    \theutterance \stepcounter{utterance}  
    & & & \multicolumn{4}{p{0.6\linewidth}}{
        \cellcolor[rgb]{0.9,0.9,0.9}{
            \makecell[{{p{\linewidth}}}]{
                \texttt{\tiny{[P1$\langle$GM]}}
                \texttt{Game Violation: ParseError: Response is not valid JSON.} \\
\texttt{Output ONLY:} \\
\texttt{One JSON with confidence\_score, current\_mode, input\_rationale, and hypothesis keys ...} \\
\texttt{OR} \\
\texttt{SOLVE: ```python ... ```} \\
            }
        }
    }
     \\ \\

    \theutterance \stepcounter{utterance}  
    & \multicolumn{4}{p{0.6\linewidth}}{
        \cellcolor[rgb]{0.561,0.737,0.561}{
            \makecell[{{p{\linewidth}}}]{
                \texttt{\tiny{[P1$\rangle$GM]}}
                \texttt{```json} \\
\texttt{\{} \\
\texttt{  "current\_mode": "probing",} \\
\texttt{  "input\_rationale": "The preloaded example f(1, 10) = 3 suggests the function counts something between a and b. Possible candidates: primes (2,3,5,7 → 4, not 3), odd numbers (1,3,5,7,9 → 5, not 3), or numbers with an odd number of set bits in binary (1:1, 2:10, 3:11, 4:100, 5:101, 6:110, 7:111, 8:1000, 9:1001, 10:1010 → 1,2,4,7,8,10 → 6, not 3). Another guess: count of numbers where a * x + b has a specific property. Alternatively, count of numbers where the digit sum is prime (1:1, 2:2, 3:3, 4:4, 5:5, 6:6, 7:7, 8:8, 9:9, 10:1 → 2,3,5,7 → 4, not 3). I'll test a hypothesis that it counts numbers where the binary representation has an odd number} \\
            }
        }
    }
    & & \\ \\

    \theutterance \stepcounter{utterance}  
    & & & \multicolumn{4}{p{0.6\linewidth}}{
        \cellcolor[rgb]{0.9,0.9,0.9}{
            \makecell[{{p{\linewidth}}}]{
                \texttt{\tiny{[P1$\langle$GM]}}
                \texttt{Game Violation: ParseError: Response is not valid JSON.} \\
\texttt{Output ONLY:} \\
\texttt{One JSON with confidence\_score, current\_mode, input\_rationale, and hypothesis keys ...} \\
\texttt{OR} \\
\texttt{SOLVE: ```python ... ```} \\
            }
        }
    }
     \\ \\

    \theutterance \stepcounter{utterance}  
    & \multicolumn{4}{p{0.6\linewidth}}{
        \cellcolor[rgb]{0.561,0.737,0.561}{
            \makecell[{{p{\linewidth}}}]{
                \texttt{\tiny{[P1$\rangle$GM]}}
                \texttt{```json} \\
\texttt{\{} \\
\texttt{  "current\_mode": "probing",} \\
\texttt{  "input\_rationale": "The preloaded example f(1, 10) = 3 suggests the function counts something between a and b. A plausible candidate is the count of numbers where the number of set bits (1s) in their binary representation is odd. For 1 to 10: 1 (1), 2 (10), 3 (11), 4 (100), 5 (101), 6 (110), 7 (111), 8 (1000), 9 (1001), 10 (1010). The counts of 1s are 1,1,2,1,2,2,3,1,2,2. Odd counts: 1,2,4,7,8 → 5 numbers, not 3. Another guess: count of numbers where the digit sum is odd (1,3,5,7,9,10 → 6, not 3). Alternatively, count of numbers where the number is a power of 2 (1,2,4,8 → 4, not 3). I'll request a new input to narrow it down.",} \\
\texttt{  "hypothesis": "unknown",} \\
\texttt{  "confidence\_score": 30} \\
\texttt{\}} \\
\texttt{```} \\
            }
        }
    }
    & & \\ \\

    \theutterance \stepcounter{utterance}  
    & & & \multicolumn{4}{p{0.6\linewidth}}{
        \cellcolor[rgb]{0.9,0.9,0.9}{
            \makecell[{{p{\linewidth}}}]{
                \texttt{\tiny{[P1$\langle$GM]}}
                \texttt{OUTPUT: ((({-}10, 1), 14), False)} \\
            }
        }
    }
     \\ \\

    \theutterance \stepcounter{utterance}  
    & \multicolumn{4}{p{0.6\linewidth}}{
        \cellcolor[rgb]{0.561,0.737,0.561}{
            \makecell[{{p{\linewidth}}}]{
                \texttt{\tiny{[P1$\rangle$GM]}}
                \texttt{```json} \\
\texttt{\{} \\
\texttt{  "current\_mode": "probing",} \\
\texttt{  "input\_rationale": "The new input f({-}10, 1) = 14 is incorrect, meaning my previous hypotheses are wrong. The negative range suggests the function might involve absolute values or counting numbers with certain properties regardless of sign. The output 14 for ({-}10, 1) is large, so it might count all integers in the range (12 numbers: {-}10 to 1 inclusive), but 14 is higher. Alternatively, it could count the sum of digits (e.g., {-}10: 1+0=1, {-}9:9, ..., 0:0, 1:1 → sum is 1+9+8+...+0+1=54, not 14). Another idea: count of numbers where the binary representation has an even number of 1s (opposite of earlier). For {-}10 to 1, this is tricky due to negative numbers' binary representation, but if we consider two's complement, it's complex. Alternatively, count of numbers where the absolute value is a prime (2,3,5,7 for 1{-}10 → 4, not 3). Not matching. I'll try a new hypothesis: count of numbers where the number of letters in the English word is odd (e.g., 'one'=3, 'two'} \\
            }
        }
    }
    & & \\ \\

    \theutterance \stepcounter{utterance}  
    & & & \multicolumn{4}{p{0.6\linewidth}}{
        \cellcolor[rgb]{0.9,0.9,0.9}{
            \makecell[{{p{\linewidth}}}]{
                \texttt{\tiny{[P1$\langle$GM]}}
                \texttt{Game Violation: ParseError: Response is not valid JSON.} \\
\texttt{Output ONLY:} \\
\texttt{One JSON with confidence\_score, current\_mode, input\_rationale, and hypothesis keys ...} \\
\texttt{OR} \\
\texttt{SOLVE: ```python ... ```} \\
            }
        }
    }
     \\ \\

    \theutterance \stepcounter{utterance}  
    & \multicolumn{4}{p{0.6\linewidth}}{
        \cellcolor[rgb]{0.561,0.737,0.561}{
            \makecell[{{p{\linewidth}}}]{
                \texttt{\tiny{[P1$\rangle$GM]}}
                \texttt{```json} \\
\texttt{\{} \\
\texttt{  "current\_mode": "probing",} \\
\texttt{  "input\_rationale": "The new input f({-}10, 1) = 14 (marked False) suggests the function is not counting primes, digit sums, or binary properties. The output 14 for ({-}10, 1) is close to the total numbers in the range (12), so it might be related to the count of numbers with a specific property that's common. Another possibility: the function counts the number of integers where the absolute value is a Fibonacci number (1, 2, 3, 5, 8, 13). For ({-}10, 1), the absolute values are 10,9,8,7,6,5,4,3,2,1,0. Fibonacci numbers here: 1,2,3,5,8 → 5 numbers, not 14. Alternatively, count of numbers where the absolute value is a perfect square (0,1,4,9). For ({-}10,1): 0,1,4,9 → 4 numbers, not 14. I'll request another input to narrow it down.",} \\
\texttt{  "hypothesis": "unknown",} \\
\texttt{  "confidence\_score": 20} \\
\texttt{\}} \\
\texttt{```} \\
            }
        }
    }
    & & \\ \\

    \theutterance \stepcounter{utterance}  
    & & & \multicolumn{4}{p{0.6\linewidth}}{
        \cellcolor[rgb]{0.9,0.9,0.9}{
            \makecell[{{p{\linewidth}}}]{
                \texttt{\tiny{[P1$\langle$GM]}}
                \texttt{OUTPUT: (((10, {-}99), 4), True)} \\
            }
        }
    }
     \\ \\

    \theutterance \stepcounter{utterance}  
    & \multicolumn{4}{p{0.6\linewidth}}{
        \cellcolor[rgb]{0.561,0.737,0.561}{
            \makecell[{{p{\linewidth}}}]{
                \texttt{\tiny{[P1$\rangle$GM]}}
                \texttt{```json} \\
\texttt{\{} \\
\texttt{  "current\_mode": "probing",} \\
\texttt{  "input\_rationale": "The new input f(10, {-}99) = 4 (marked True) is very revealing. The function seems to work regardless of argument order (since 10 > {-}99 but returns 4, same as if we'd passed ({-}99, 10)). The output 4 suggests it's counting something specific in the range. Possible candidates: count of single{-}digit numbers (10 is not single{-}digit, {-}99 has digits 9 and 9 → 0, but output is 4, so no). Count of numbers where the absolute value is a single{-}digit prime (2,3,5,7). For (10, {-}99), absolute values are 10 to 99 → primes: 11,13,...,97 → 21 primes, not 4. Alternatively, count of numbers where the absolute value is a power of 2 (1,2,4,8,16,32,64). For (10, {-}99): 16,32,64 → 3 numbers, not 4. Another idea: count of numbers where the digit sum is 1 (1,10,{-}9 → absolute values: 1,10,9 → only 1 has digit sum 1). Not matching.} \\
            }
        }
    }
    & & \\ \\

    \theutterance \stepcounter{utterance}  
    & & & \multicolumn{4}{p{0.6\linewidth}}{
        \cellcolor[rgb]{0.9,0.9,0.9}{
            \makecell[{{p{\linewidth}}}]{
                \texttt{\tiny{[P1$\langle$GM]}}
                \texttt{Game Violation: ParseError: Response is not valid JSON.} \\
\texttt{Output ONLY:} \\
\texttt{One JSON with confidence\_score, current\_mode, input\_rationale, and hypothesis keys ...} \\
\texttt{OR} \\
\texttt{SOLVE: ```python ... ```} \\
            }
        }
    }
     \\ \\

    \theutterance \stepcounter{utterance}  
    & \multicolumn{4}{p{0.6\linewidth}}{
        \cellcolor[rgb]{0.561,0.737,0.561}{
            \makecell[{{p{\linewidth}}}]{
                \texttt{\tiny{[P1$\rangle$GM]}}
                \texttt{```json} \\
\texttt{\{} \\
\texttt{  "current\_mode": "probing",} \\
\texttt{  "input\_rationale": "The input f(10, {-}99) = 4 (True) suggests the function counts something consistent regardless of argument order. The output 4 is small, so it likely counts numbers with a very specific property. Possible candidates: count of numbers where the absolute value is a perfect cube (1, 8, 27, 64). For (10, {-}99), absolute values are 10 to 99 → cubes: 27, 64 → 2 numbers, not 4. Alternatively, count of numbers where the absolute value is a perfect square (1, 4, 9, 16, 25, 36, 49, 64, 81). For (10, {-}99): 16, 25, 36, 49, 64, 81 → 6 numbers, not 4. Another idea: count of numbers where the absolute value is a Fibonacci number (1, 2, 3, 5, 8, 13, 21, 34, 55, 89). For (10, {-}99): 13, 21, 34, 55, 89 →} \\
            }
        }
    }
    & & \\ \\

    \theutterance \stepcounter{utterance}  
    & & & \multicolumn{4}{p{0.6\linewidth}}{
        \cellcolor[rgb]{0.9,0.9,0.9}{
            \makecell[{{p{\linewidth}}}]{
                \texttt{\tiny{[P1$\langle$GM]}}
                \texttt{Game Violation: ParseError: Response is not valid JSON.} \\
\texttt{Output ONLY:} \\
\texttt{One JSON with confidence\_score, current\_mode, input\_rationale, and hypothesis keys ...} \\
\texttt{OR} \\
\texttt{SOLVE: ```python ... ```} \\
            }
        }
    }
     \\ \\

    \theutterance \stepcounter{utterance}  
    & \multicolumn{4}{p{0.6\linewidth}}{
        \cellcolor[rgb]{0.561,0.737,0.561}{
            \makecell[{{p{\linewidth}}}]{
                \texttt{\tiny{[P1$\rangle$GM]}}
                \texttt{```json} \\
\texttt{\{} \\
\texttt{  "current\_mode": "probing",} \\
\texttt{  "input\_rationale": "Combining all examples: f(1,10)=3 (True), f({-}10,1)=14 (False), f(10,{-}99)=4 (True). The function seems to count numbers where the absolute value is a single{-}digit number (0{-}9). For (1,10): 1{-}9 → 9 numbers, but output is 3 (no). For (10,{-}99): {-}9 to 9 → 19 numbers, but output is 4 (no). Another idea: count of numbers where the absolute value is a prime power (primes or powers of primes). For (1,10): 2,3,4,5,7,8,9 → 7 numbers, not 3. For (10,{-}99): too many, not 4. Final guess: count of numbers where the absolute value is a factorial (1,2,6,24). For (1,10): 1,2,6 → 3 numbers (matches f(1,10)=3). For (10,{-}99): 24 → 1 number (but output is 4, so no).",} \\
\texttt{  "hypothesis": "def f(a: int, b: int) {-}> int:\textbackslash{}n    \textbackslash{}"\textbackslash{}"} \\
            }
        }
    }
    & & \\ \\

    \theutterance \stepcounter{utterance}  
    & & & \multicolumn{4}{p{0.6\linewidth}}{
        \cellcolor[rgb]{0.9,0.9,0.9}{
            \makecell[{{p{\linewidth}}}]{
                \texttt{\tiny{[P1$\langle$GM]}}
                \texttt{Game Violation: ParseError: Response is not valid JSON.} \\
\texttt{Output ONLY:} \\
\texttt{One JSON with confidence\_score, current\_mode, input\_rationale, and hypothesis keys ...} \\
\texttt{OR} \\
\texttt{SOLVE: ```python ... ```} \\
            }
        }
    }
     \\ \\

    \theutterance \stepcounter{utterance}  
    & \multicolumn{4}{p{0.6\linewidth}}{
        \cellcolor[rgb]{0.561,0.737,0.561}{
            \makecell[{{p{\linewidth}}}]{
                \texttt{\tiny{[P1$\rangle$GM]}}
                \texttt{```json} \\
\texttt{\{} \\
\texttt{  "current\_mode": "confirming",} \\
\texttt{  "input\_rationale": "After analyzing all examples, the function appears to count numbers where the absolute value is a prime number less than 10 (2, 3, 5, 7). For (1,10): 2,3,5,7 → 4 primes, but output is 3 (mismatch). However, if we exclude 2 (only even prime), it matches f(1,10)=3 (3,5,7). For (10,{-}99): 3,5,7 → 3 primes, but output is 4 (mismatch). Final hypothesis: count of numbers where the absolute value is a Fibonacci number (1,2,3,5,8). For (1,10): 1,2,3,5,8 → 5 numbers, but output is 3. If we exclude 1 and 8: 2,3,5 → 3 (matches f(1,10)=3). For (10,{-}99): 13,21,34,55,89 → 5 numbers, but output is 4. If we exclude 89: 13,21,34,55 → 4 (matches f(10,{-}99)=4). This seems consistent.",} \\
            }
        }
    }
    & & \\ \\

    \theutterance \stepcounter{utterance}  
    & & & \multicolumn{4}{p{0.6\linewidth}}{
        \cellcolor[rgb]{0.9,0.9,0.9}{
            \makecell[{{p{\linewidth}}}]{
                \texttt{\tiny{[P1$\langle$GM]}}
                \texttt{Game Violation: ParseError: Response is not valid JSON.} \\
\texttt{Output ONLY:} \\
\texttt{One JSON with confidence\_score, current\_mode, input\_rationale, and hypothesis keys ...} \\
\texttt{OR} \\
\texttt{SOLVE: ```python ... ```} \\
            }
        }
    }
     \\ \\

    \theutterance \stepcounter{utterance}  
    & \multicolumn{4}{p{0.6\linewidth}}{
        \cellcolor[rgb]{0.561,0.737,0.561}{
            \makecell[{{p{\linewidth}}}]{
                \texttt{\tiny{[P1$\rangle$GM]}}
                \texttt{```json} \\
\texttt{\{} \\
\texttt{  "current\_mode": "confirming",} \\
\texttt{  "input\_rationale": "Based on the examples, the function appears to count numbers in the range {[}min(a,b), max(a,b){]} where the absolute value is a Fibonacci number (1, 2, 3, 5, 8, 13, 21, 34, 55, 89), excluding 1 and 89. This matches: f(1,10)=3 (2,3,5), f(10,{-}99)=4 (13,21,34,55). The preloaded example f({-}10,1)=14 was marked False, suggesting my hypothesis is incomplete. Final attempt: count all Fibonacci numbers in range, including 1 and 89.",} \\
\texttt{  "hypothesis": "def f(a: int, b: int) {-}> int:\textbackslash{}n    \textbackslash{}"\textbackslash{}"\textbackslash{}"Count numbers in {[}a,b{]} where absolute value is a Fibonacci number.\textbackslash{}"\textbackslash{}"\textbackslash{}"\textbackslash{}n    fibs = \{1, 2, 3, 5, 8, 13, 21, 34, 55, 89\}\textbackslash{}n    low, high = min(a, b), max(a, b)\textbackslash{}n    return sum(1 for x in range(low, high + 1) if abs(x)} \\
            }
        }
    }
    & & \\ \\

    \theutterance \stepcounter{utterance}  
    & & & \multicolumn{4}{p{0.6\linewidth}}{
        \cellcolor[rgb]{0.9,0.9,0.9}{
            \makecell[{{p{\linewidth}}}]{
                \texttt{\tiny{[P1$\langle$GM]}}
                \texttt{Game Violation: ParseError: Response is not valid JSON.} \\
\texttt{Output ONLY:} \\
\texttt{One JSON with confidence\_score, current\_mode, input\_rationale, and hypothesis keys ...} \\
\texttt{OR} \\
\texttt{SOLVE: ```python ... ```} \\
            }
        }
    }
     \\ \\

    \theutterance \stepcounter{utterance}  
    & \multicolumn{4}{p{0.6\linewidth}}{
        \cellcolor[rgb]{0.561,0.737,0.561}{
            \makecell[{{p{\linewidth}}}]{
                \texttt{\tiny{[P1$\rangle$GM]}}
                \texttt{SOLVE:} \\
\texttt{```python} \\
\texttt{def f(a: int, b: int) {-}> int:} \\
\texttt{    """Count numbers in {[}a,b{]} where the absolute value is a Fibonacci number (1, 2, 3, 5, 8, 13, 21, 34, 55, 89)."""} \\
\texttt{    fibs = \{1, 2, 3, 5, 8, 13, 21, 34, 55, 89\}} \\
\texttt{    low, high = min(a, b), max(a, b)} \\
\texttt{    return sum(1 for x in range(low, high + 1) if abs(x) in fibs)} \\
\texttt{```} \\
            }
        }
    }
    & & \\ \\

    \theutterance \stepcounter{utterance}  
    & & & \multicolumn{2}{p{0.3\linewidth}}{
        \cellcolor[rgb]{0.9,0.9,0.9}{
            \makecell[{{p{\linewidth}}}]{
                \texttt{\tiny{[GM$|$GM]}}
                \texttt{Game Verdict: LOSS} \\
            }
        }
    }
    & & \\ \\

    \theutterance \stepcounter{utterance}  
    & & & \multicolumn{2}{p{0.3\linewidth}}{
        \cellcolor[rgb]{0.9,0.9,0.9}{
            \makecell[{{p{\linewidth}}}]{
                \texttt{\tiny{[GM$|$GM]}}
                \texttt{Incorrect.} \\
\texttt{Failed Cases:} \\
\texttt{Input: {[}1, 10{]} | Expected: 3 | Got: 5} \\
\texttt{Input: {[}{-}10, 1{]} | Expected: 4 | Got: 6} \\
\texttt{Input: {[}10, {-}99{]} | Expected: 4 | Got: 15} \\
            }
        }
    }
    & & \\ \\

    \theutterance \stepcounter{utterance}  
    & & & \multicolumn{2}{p{0.3\linewidth}}{
        \cellcolor[rgb]{0.9,0.9,0.9}{
            \makecell[{{p{\linewidth}}}]{
                \texttt{\tiny{[GM$|$GM]}}
                \texttt{def product\_mod\_seven(a: int, b: int) {-}> int:} \\
\texttt{    return (a * b) \% 7} \\
            }
        }
    }
    & & \\ \\

    \theutterance \stepcounter{utterance}  
    & & & \multicolumn{2}{p{0.3\linewidth}}{
        \cellcolor[rgb]{0.9,0.9,0.9}{
            \makecell[{{p{\linewidth}}}]{
                \texttt{\tiny{[GM$|$GM]}}
                \texttt{def f(a: int, b: int) {-}> int:} \\
\texttt{    """Count numbers in {[}a,b{]} where the absolute value is a Fibonacci number (1, 2, 3, 5, 8, 13, 21, 34, 55, 89)."""} \\
\texttt{    fibs = \{1, 2, 3, 5, 8, 13, 21, 34, 55, 89\}} \\
\texttt{    low, high = min(a, b), max(a, b)} \\
\texttt{    return sum(1 for x in range(low, high + 1) if abs(x) in fibs)} \\
            }
        }
    }
    & & \\ \\

    \theutterance \stepcounter{utterance}  
    & & & \multicolumn{2}{p{0.3\linewidth}}{
        \cellcolor[rgb]{0.9,0.9,0.9}{
            \makecell[{{p{\linewidth}}}]{
                \texttt{\tiny{[GM$|$GM]}}
                \texttt{Mode: passive\_labeled\_pairs} \\
\texttt{Accuracy: 0.100} \\
\texttt{Information requests: 2/15} \\
\texttt{Revealed passive examples: 3} \\
\texttt{Efficiency: 0.933} \\
\texttt{Slack: 13} \\
\texttt{Observed{-}pair consistency: 0.000} \\
\texttt{Membership balance: 2 true / 1 false} \\
\texttt{Parse errors: 8} \\
\texttt{Runtime errors: 0} \\
            }
        }
    }
    & & \\ \\

\end{supertabular}
}

\end{document}
